\chapter{Ce que dit vraiment la dopamine}
% version complète: chapters/07_dofaalt.tex

Au chapitre précédent, la dopamine était simple: un seul nombre, \og mieux ou moins bien que prévu\fg{}. C'est une première approximation, et elle marche bien. Mais les expériences de ces dernières années ont montré que ce \og fil\fg{} transporte plus qu'un nombre. Ce chapitre parle des endroits où le tableau se complique et où les chercheurs débattent encore.

\section*{Optimistes et pessimistes}

Les neurones à dopamine ne sont pas tous pareils. Certains réagissent plus fort aux bonnes surprises qu'aux mauvaises: ce sont les optimistes. D'autres font l'inverse: les pessimistes. Si chaque neurone apprend selon sa propre règle, l'optimiste apprendra ce que sera la récompense dans le meilleur des cas, et le pessimiste dans le pire. Ensemble, ils n'enregistrent pas une seule attente moyenne, mais tout l'éventail des issues possibles, comme un bookmaker qui connaît non seulement le favori, mais aussi la cote de chacun. Cette idée s'accorde avec l'expérience; à quoi sert au cerveau tout l'éventail, c'est encore une hypothèse.

\pencilfig{7}{\pencilart{7}}{Chaque neurone à dopamine tient son propre point sur la courbe des récompenses possibles: ensemble, ils mémorisent toute la courbe.}

\section*{Pas seulement \og mieux\fg{}, mais aussi \og important\fg{}}

Une partie des neurones à dopamine s'embrase aussi pour des événements désagréables mais importants: un bruit fort, un danger. Il semble que le réseau dopaminergique transporte non seulement le signe de la surprise, mais aussi son importance, un signal \og fais attention\fg{}.

\section*{Deux signaux sur un même fil}

Les bouffées rapides font apprendre. Mais le niveau moyen de dopamine, lent, qui varie en quelques minutes, semble dire autre chose: à quel point l'environnement est riche en ce moment. Si les récompenses abondent alentour, le temps est précieux, et l'animal agit plus vite. Les expériences confirment que le niveau de dopamine est lié à l'énergie des actions. Un ingénieur appellerait cela un partage des signaux par fréquence: les hautes fréquences portent une chose, les basses une autre.

\section*{La montée près du but}

Quand l'animal s'approche d'une récompense, la dopamine souvent ne s'embrase pas, mais monte doucement. Première explication: c'est l'attente elle-même, un signal \og le but est proche, accélère\fg{}. Une autre garde l'idée de surprise: de loin, l'animal évalue la situation de façon floue, puis de plus en plus précisément en approchant, et chaque précision est une petite bonne surprise. Une élégante expérience dans un couloir virtuel, où l'animal était soudain téléporté plus près du but, plaide pour la seconde: la dopamine répondait par une bouffée, comme il se doit pour une surprise.

\section*{Le regard en arrière}

Il existe aussi une alternative audacieuse: la dopamine ne répondrait pas à la question \og que va-t-il se passer?\fg{}, mais à \og qu'est-ce qui a causé la récompense?\fg{}; elle regarderait vers le passé, non vers l'avenir. Les expériences où les deux théories prédisent des choses différentes se contredisent encore. C'est un débat scientifique bien vivant.

\section*{Un vecteur au lieu d'un nombre}

Enfin, la dopamine réagit aussi au changement de \emph{goût} de la récompense quand sa valeur n'a pas changé, et elle aide à apprendre sans aucune récompense. On dirait non pas une seule erreur, mais tout un faisceau d'erreurs, une pour chaque trait de ce qui s'est passé.

Conclusion du chapitre: la plupart des nouvelles théories n'annulent pas le tableau simple, elles l'élargissent. La dopamine n'est pas un fil unique, mais un petit faisceau de quelques fils. Seul le \og regard en arrière\fg{} le conteste sérieusement.

\fullversion{7}
